\scanpage{92}
\mspage{91}
opérer $\mu$ trivialement sur $C$ et sur $\Pi_1(T)$.
On va relier les op. de $\mu$ sur $\Pi_1(P)$
resp. sur $C'$ avec la $\pi$-structure sur
$\Pi_1(p)$, ou celle de $f$ qui s'en déduit
par transport de structure. Considérons
pour ceci la suite exacte
\[
(12)\qquad 0\longrightarrow\pi\longrightarrow\widetilde U\longrightarrow U\longrightarrow0
\]
et son image inverse par (6), d'où
\[
(13)\qquad 0\mathbin{\text{---}}\pi\longrightarrow\mathcal E\longrightarrow\mu\longrightarrow0
\]
\margin{ext. \unsure{commutative}}
\deleted{On suppose}
\deleted{Supposons que $\mu$}
Soient $x,y$ deux points de $\unsure{E}$ \uline{qui soient}
\uline{congrus mod} $\mu$,
\inserted{i.e. dans une \deleted{\unsure{la}} même fibre $C'_{\unsure{\xi}}$ de $C'$ \unsure{sur $C$}}
(\deleted{\unsure{dont}} On suppose que \deleted{\unsure{il}} \unsure{$\mu$} y
opère librement sur $E$ (automatique si $i_E$ injectif)).
Je dis qu'on a alors un isom canonique
\inserted{de $\pi$-torseurs}
\[
(14)\qquad \operatorname{Isom}_{C'_{\unsure{\xi}}}(x,y)\simeq\mathcal E_{y-x}
\]
où $\unsure{z}=y-x$ est l'unique $z\in\mu$ tel que
$y=\unsure{z}\mathbin{.}x$, et où $\mathcal E_{\unsure{z}}$ désigne la fibre de
l'ext. $\mathcal E$ (13) sur l'élément $z\in\mu$. Cet
isom (14) est compatible avec la composition
dans $C'_{\unsure{\xi}}$, qui correspond à l'addition de $\mathcal E$.
